\documentclass[book.tex]{subfiles}

\begin{document}
\chapter{General Relativity}

\label{chap:gr}
\noindent\rule{11cm}{0.4pt} \\
\textbf{Prerequisites:} \autoref{chap:vectors}, \autoref{chap:manifold}, \autoref{chap:special_relativity}  \\
\textbf{Difficulty Level:} *** \\
\noindent\rule{11cm}{0.4pt}
\vspace{5mm}

General relativity is the study of accelerating relativistic systems. Unlike special relativity (which we studied in Chapter~\ref{chap:special_relativity}), general relativistic systems require more mathematical foundations from manifolds. In particular, it becomes necessary to deal with tensors and curvature. A lot of the math involved in general relativity is different from that used elsewhere in physics which creates additional complications \cite{carroll2001no}.

Note that when we discuss tensors in this chapter, we are actually talking about tensor fields on a spacetime. This difference can prove a source of considerable confusion for newcomers who mistakenly identify "tensors" in machine learning (multidimensional arrays) with the tensor fields of physics.

\section{Reviewing Special Relativity}

We've previously introduced the basics of special relativity, but we will review it here as well. Special relativity occurs on a special four dimensional manifold, the Minkowski spacetime. By convention, the coordinates of a point in Minkowski space time are written

\begin{align}
    x^{\mu} = (t, x, y, z)
\end{align}

Here, we assume as we have done earlier that the speed of light $c = 1$. The elements of the Minkowski spacetime are called events. We also follow the convention that an event in space time is denoted with a capital letter $V^{\mu}$. (Note the difference with the usual convention in which elements of a vector space are denoted with lower case elements). Minkowski spacetime differs from the usual space of $\mathbb{R}^4$ in one critical fashion: distances are computed differently. In particular we define

\begin{equation}
    \eta_{\mu\nu} = \begin{pmatrix}
    -1 & 0 & 0 & 0 \\
    0  & 1 & 0 & 0 \\
    0  & 0 & 1 & 0 \\
    0  & 0 & 0 & 1
    \end{pmatrix}
\end{equation}

We can then define the inner product between two events $A$ and $B$ as

\begin{align}
    A \cdot B &\equiv \eta_{\mu\nu} A^{\mu} B^{\nu} \\
    &= - A^0B^0 + A^1 B^1 + A^2 B^2 + A^3 B^3
\end{align}

Here we are following tensor summation convention that indices mentioned in both sub and superscripts are implicitly summed over. This convention lends itself well to describing the distance between two points as 
\begin{align}
    ds^2 &= \eta_{\mu\nu} dx^{\mu} dx^{\nu} \\
    &= -dt^2 + dx^2 + dy^2 + dz^2
\end{align}
This is the metric associated with Minkowski space. Note in particular that distances can be negative. The distance between $(t, x, y, z)$ and $(2t, x, y, z)$ is negative for example. We consequently define the proper time infinitesimally as
\begin{align}
    d\tau^2 = - ds^2
\end{align}
and define the proper time $\tau$ as
\begin{align}
    \tau = \int \sqrt{-ds^2}
\end{align}


The choice of metric also suggests a few terms. A vector $V^{\mu}$ is timelike if $V \cdot V < 0$. (Note why this makes sense.) A vector $W^{\mu}$ such that $W \cdot W > 0$ is spacelike, and a vector with 0 norm is called null. A path through spacetime is denoted $x^{\mu}(\lambda)$. We define the four-velocity as 
\begin{equation}
    U^{\mu} = \frac{dx^{\mu}}{d\tau}
\end{equation}
The definition of the momentum four-vector follows as
\begin{equation}
    p^{\mu} = m U^{\mu}
\end{equation}
where $m$ is the mass of a particle. 
% Check out \url{https://preposterousuniverse.com/wp-content/uploads/2015/08/grtinypdf.pdf}



\section{General Relativity as a Physical Theory}

%\textcolor{red}{TODO: Define the Ricci tensor}\\
%\textcolor{red}{TODO: Define the cosmological constant}\\
%\textcolor{red}{TODO: Define the energy momentum tensor}

A pseudo-Riemannian manifold $(M, g)$ is a differentiable manifold $M$ equipped with a smooth, symmetric metric tensor $g$. Spacetime $\mathbb{R}^{1,3}$ with the usual metric
\begin{align}
    -dt^2 + dx^2 + dy^2 + dz^2
\end{align}
is a pseudo-Riemannian manifold.


The Einstein tensor is defined as 
\begin{align}
    G_{\mu \nu} = R_{\mu \nu} - \frac{1}{2} R g_{\mu \nu} 
\end{align}
where $R_{\mu \nu}$ is the Ricci tensor, $R$ the scalar curvature, and $g_{\mu \nu}$ is the metric tensor. The Ricci tensor is defined in terms of the Riemann tensor
\begin{align}
R_{ab} &= R^{c}_{abc}
\end{align}
The Riemann tensor itself is expanded in terms of Christoffel symbols
\begin{align}
R^{\rho}_{\sigma \mu \nu} &= \partial_\mu \Gamma^\rho_{\nu \sigma} - \partial \nu \Gamma^\rho_{\mu \sigma} + \Gamma^\rho_{\mu \lambda} \Gamma^\lambda_{\nu \sigma} - \Gamma^\rho_{\nu \lambda} \Gamma^\lambda_{\mu \sigma}
\end{align}

The Einstein gravitational constant is defined as
\begin{align}
    \kappa = \frac{8 \pi G}{c^4}
\end{align}
where $G$ is the Newtonian gravitational constant, and $c$ is the speed of light in a vacuum.

%\textcolor{red}{TODO: Introduce einstein field equations}

\begin{align}
    G_{ab} &= R_{ab} - \frac{R}{2} g_{ab} = 8 \pi T_{ab}
\end{align}
where $G_{ab}$ is the Einstein tensor, $R_{ab}$ is the Ricci tensor of the spacetime, $R$ is the Ricci scalar, and $T_{ab}$ is the stress-energy tensor. The Ricci tensor can be expanded in terms of Christophel symbols as follows:

\begin{align}
    R_{ab} &= \partial_c \Gamma_{ab}^c - \partial_a \Gamma^c_{cb} + \Gamma^c_{cd}\Gamma^d_{ab} - \Gamma^c_{da}\Gamma^d_{cb} \\
    \Gamma^c_{ab} &= \frac{1}{2}g_{ab}(\partial_a g_{bd} + \partial_b g_{ad} - \partial_d g_{ab})
\end{align}

%\textcolor{red}{TODO: What is the intuition behind these definitions? Is there a natural geometric/mathematical explanation we can point to? We mention Christophel symbols previously in the preliminaries. What's the chapter reference?}

The equations above define a set of 10 coupled partial differential equations %(TODO: why 10? prove in more detail)

\section{General Relativity as a Physical Theory}

Let us now formulate general relativity as a physical theory in our framework. The state space $\mathcal{S}$ for general relativity is the category of pseudo-Riemannian manifolds %\textcolor{red}{TODO: Is this right? It feels too big}.

Time evolution in general relativity is a bit stranger than it is for other equations. A solution describes a configuration of spacetime in which there isn't an explicit notion of time evolution. The equations to be solved of course are the Einstein field equation
\begin{align}
    G_{\mu \nu} = \Lambda g_{\mu \nu} + \kappa T_{\mu \nu}
\end{align}
Here $G_{\mu \nu} $ is the Einstein tensor.

Parameters are given by
\begin{align}\mathcal{W} = {G, c}\end{align} 
where $G$ is the Newtonian gravity constant, and $c$ is the speed of light in a vacuum. Domain of Applicability $\mathcal{C}$ constrains general relativity to only be valid for suitably macroscopic systems (not quantum scale).

%\textcolor{red}{TODO: What are the right hyperparameters for this system?} 

%\textcolor{red}{TODO: What are measurable observables in general relativity? Perhaps tie to some of the classical experiments?}

%\section{Schwarzschild Metric}

\section{Exercises}
\begin{enumerate}
    \item Prove that the four-velocity satisfies the equation
    \begin{equation}
        \eta_{\mu\nu} U^{\mu} U^{\nu} = -1
    \end{equation}
\end{enumerate}
\end{document}
